\documentclass[a4paper,12pt]{article}
\usepackage[T2A]{fontenc}
\usepackage[utf8]{inputenc}
\usepackage[russian]{babel}
\usepackage{amsmath,amssymb,mathtools}
\renewcommand{\familydefault}{\sfdefault}
\usepackage{icomma}
\usepackage{geometry}
\geometry{left=22mm,right=22mm,top=22mm,bottom=23mm}
\usepackage{microtype}
\usepackage{tikz}
\usetikzlibrary{arrows.meta,positioning,shapes.geometric,calc}
\usepackage{tabularx,booktabs,longtable}
\usepackage{enumitem}
\usepackage{needspace}
\usepackage{hyperref}
\usepackage{parskip}
\usepackage{setspace}
\usepackage{xcolor}
\usepackage{fancyhdr}
\definecolor{accent}{HTML}{087F82}
\definecolor{darkaccent}{HTML}{07575D}
\definecolor{textgray}{HTML}{333333}
\definecolor{muted}{HTML}{68767B}
\hypersetup{colorlinks=true,linkcolor=darkaccent,urlcolor=darkaccent}
\setlength{\parindent}{0pt}
\setlength{\parskip}{5pt}
\setstretch{1.22}
\setlist[enumerate]{leftmargin=*,itemsep=7pt,topsep=5pt}
\setlist[itemize]{leftmargin=*,itemsep=3pt,topsep=3pt}
\pagestyle{fancy}
\fancyhf{}
\fancyhead[L]{\footnotesize\color{muted} Движение в одном направлении}
\fancyhead[R]{\footnotesize\color{muted} Ярослав Гаврилов}
\fancyfoot[C]{\footnotesize\color{muted}\thepage}
\renewcommand{\headrulewidth}{0.25pt}
\setlength{\headheight}{15pt}
\newcommand{\key}[1]{\textcolor{darkaccent}{\textbf{#1}}}
\newcommand{\task}[2]{\Needspace{9\baselineskip}\par\noindent\textbf{\textcolor{darkaccent}{#1.} #2}\par}
\newcommand{\smallnote}[1]{{\small\color{muted}#1}}
\begin{document}
\thispagestyle{empty}
\begin{center}
\vspace*{27mm}
{\Large\color{darkaccent} МАТЕМАТИКА}\par
\vspace{18mm}
{\Huge\color{textgray} Чек-лист}\par
\vspace{14mm}
{\LARGE\color{accent} Текстовые задачи на движение}\par
\vspace{3mm}
{\Large\color{accent} Догонка и движение по этапам}\par
\vspace{15mm}
{\large Подготовка к экзамену по математике}\par
\vfill
{\large Ярослав Гаврилов}\par
\vspace{10mm}
{\large 12 ноября 2026 г.}\par
\vspace{15mm}
{\normalsize Лёвин Артём Александрович}\par
{\normalsize эксклюзивно для Ярослава Гаврилова}\par
\vspace{13mm}
\begin{tikzpicture}
\draw[accent,line width=0.9pt] (0,0) .. controls (1.3,.8) and (2.8,-.8) .. (4,0) .. controls (5.4,.8) and (6.8,-.8) .. (8.3,0);
\fill[darkaccent] (4,0) circle (1.6pt);
\end{tikzpicture}
\end{center}
\clearpage

\section*{1. Главный принцип: описать движение математически}

В задачах на движение важно не угадывать уравнение, а последовательно перевести текст в \key{схему}, \key{таблицу} и затем в \key{уравнение}. Для равномерного движения без остановок используются три величины:
\[
\boxed{S=vt},\qquad \boxed{t=\frac{S}{v}},\qquad \boxed{v=\frac{S}{t}}.
\]
Здесь $S$~--- расстояние, $v$~--- скорость, $t$~--- время. Если скорость дана в км/ч, расстояние выражаем в километрах, а время --- в часах.

\subsection*{Правильная последовательность}
\begin{enumerate}
\item Выделите участников движения, точки старта, направление и возможные моменты встречи, разворота и финиша.
\item Нарисуйте схему и подпишите известные расстояния. Если схема ничего не добавляет (например, оба проходят один и тот же путь), достаточно таблицы.
\item Составьте таблицу: сначала впишите то, что известно и что требуется найти. Остальные клетки заполните по формулам движения.
\item Выберите величины, которые можно сравнить: \emph{пути} или \emph{времена}. Составьте равенство, исходя из условия и схемы.
\item Проверьте единицы измерения, ограничения и смысл результата.
\end{enumerate}

\subsection*{Таблица «скорость~--- время~--- расстояние»}
\begin{center}
\renewcommand{\arraystretch}{1.45}
\begin{tabularx}{\linewidth}{@{}lXXX@{}}
\toprule
Участник & Скорость $v$ & Время $t$ & Расстояние $S$\\
\midrule
Быстрый & $u+\Delta v$ & $x$ & $x(u+\Delta v)$\\
Медленный & $u$ & $x$ & $xu$\\
\bottomrule
\end{tabularx}
\end{center}
Сначала можно обозначить скорость медленного участника через $u$, а искомое время --- через $x$. \key{Последний незаполненный столбец} вычисляется по формуле из двух предыдущих. Равные времена в таблице допустимы только при \emph{одновременном начале рассматриваемого этапа}.

\section*{2. Догонка при одновременном старте}

\textbf{Ситуация.} Быстрый автомобиль выехал из пункта $A$, грузовик одновременно выехал из пункта $B$, который на $125$ км впереди $A$. Автомобиль движется на $25$ км/ч быстрее. Нужно узнать время догонки.

\begin{center}
\begin{tikzpicture}[>=Latex]
\draw[very thick] (0,0)--(12,0);
\foreach \x/\lab in {0/A,4/B,10/\text{догонка}} {\fill[darkaccent] (\x,0) circle (1.6pt);\node[below=4pt] at (\x,0) {\lab};}
\draw[<->,accent] (0,.55)--node[above]{$125\ \text{км}$}(4,.55);
\draw[->,line width=1pt,accent] (0,-.55)--(9.7,-.55);
\draw[->,line width=1pt,muted] (4,1.05)--(9.7,1.05);
\node[left] at (0,-.55) {\small легковой};
\node[left] at (4,1.05) {\small грузовик};
\end{tikzpicture}
\end{center}

Если скорость грузовика $u$, то скорость автомобиля $u+25$. За $x$ часов автомобиль проедет $x(u+25)$ километров, а грузовик --- $xu$ километров. Путь автомобиля на $125$ км больше:
\[
x(u+25)=xu+125 \quad\Longrightarrow\quad 25x=125 \quad\Longrightarrow\quad \boxed{x=5\ \text{ч}}.
\]
\key{Почему неизвестная скорость исчезла?} Оба ехали одно и то же время, поэтому одинаковый фрагмент $xu$ сокращается. Остаётся только разница скоростей. В общем случае
\[
\boxed{\Delta S=(v_{\text{быстр}}-v_{\text{медл}})t},\qquad v_{\text{быстр}}>v_{\text{медл}}.
\]

\section*{3. Один старт: расстояние между участниками}

Если два пешехода начали движение одновременно из одной точки в одном направлении, а первый идёт на $2$ км/ч быстрее, расстояние между ними за $x$ часов станет $2x$ километров. При расстоянии $300$ м сначала переводим: $300$ м $=0{,}3$ км. Тогда
\[
(u+2)x-ux=0{,}3\quad\Longrightarrow\quad 2x=0{,}3
\quad\Longrightarrow\quad x=0{,}15\ \text{ч}=\boxed{9\ \text{мин}}.
\]
\key{Проверка:} $0{,}15\cdot60=9$. Не смешивайте метры и километры в одном уравнении.

\clearpage
\section*{4. Разворот: сначала в одну сторону, затем обратно}

\textbf{Ситуация.} Два человека одновременно вышли из одной точки к опушке леса в $6$ км. Первый идёт со скоростью $4{,}5$ км/ч, второй --- со скоростью $5{,}5$ км/ч. Быстрый дошёл до опушки и сразу пошёл обратно. На каком расстоянии $x$ от опушки они встретятся?

\begin{center}
\begin{tikzpicture}[>=Latex]
\draw[very thick] (0,0)--(12,0);
\fill[darkaccent] (0,0) circle (1.7pt); \fill[darkaccent] (12,0) circle (1.7pt);
\fill[darkaccent] (10.8,0) circle (1.7pt);
\node[below] at (0,0) {старт}; \node[below] at (12,0) {опушка};
\node[below] at (10.8,0) {встреча};
\draw[<->,accent] (0,.7)--node[above]{$6\ \text{км}$}(12,.7);
\draw[->,muted,line width=1pt] (0,-.5)--(10.7,-.5);
\draw[->,accent,line width=1pt] (0,1.3)--(11.9,1.3);
\draw[->,accent,line width=1pt] (11.9,1.3)--(10.8,1.3);
\node[above] at (11.4,1.3) {\small $x$};
\end{tikzpicture}
\end{center}

Первый проходит $6-x$ км, второй --- $6+x$ км. Их \emph{общее время от старта до встречи} одинаково:
\[
\frac{6-x}{4{,}5}=\frac{6+x}{5{,}5}.
\]
Умножаем обе части на $4{,}5\cdot5{,}5$:
\[
5{,}5(6-x)=4{,}5(6+x)\quad\Longrightarrow\quad 33-5{,}5x=27+4{,}5x,
\]
\[
10x=6\quad\Longrightarrow\quad \boxed{x=0{,}6\ \text{км}=600\ \text{м}}.
\]
\key{Идея:} быстрый участник успел пройти \emph{до разворота} и \emph{после разворота}. Поэтому его путь равен $6+x$, а не $6-x$. Ограничение: $0<x<6$.

\clearpage
\section*{5. Старт в разное время и общий финиш}

\textbf{Ситуация.} Грузовик выехал из $A$ в $B$ на $6$ минут раньше автобуса. Расстояние составляет $24$ км. Скорость автобуса на $12$ км/ч больше, оба прибыли одновременно. Найдите скорость автобуса.

Обозначим скорость автобуса через $x$ км/ч. Тогда скорость грузовика $x-12$ км/ч, причём $x>12$. Оба преодолели $24$ км:
\begin{center}
\renewcommand{\arraystretch}{1.35}
\begin{tabularx}{\linewidth}{@{}lXXX@{}}
\toprule
Участник & $v$, км/ч & $t$, ч & $S$, км\\
\midrule
Автобус & $x$ & $\dfrac{24}{x}$ & $24$\\
Грузовик & $x-12$ & $\dfrac{24}{x-12}$ & $24$\\
\bottomrule
\end{tabularx}
\end{center}

Грузовик был в пути на $6$ минут больше, то есть на $\dfrac{6}{60}=\dfrac1{10}$ часа:
\[
\frac{24}{x-12}-\frac{24}{x}=\frac1{10}.
\]
После приведения к общему знаменателю $x(x-12)=2880$, откуда $x^2-12x-2880=0$. Корни $60$ и $-48$; второй не подходит, поскольку $x>12$. Значит, $\boxed{x=60\ \text{км/ч}}$.

\key{Проверка:} автобус ехал $24/60=0{,}4$ ч, грузовик --- $24/48=0{,}5$ ч. Разница $0{,}1$ ч, или $6$ минут.

\clearpage
\section*{6. Два этапа движения: две таблицы и два уравнения}

Если после встречи кто-то разворачивается либо участники продолжают путь по-разному, разделите сюжет на \emph{этап до встречи} и \emph{этап после встречи}. Для каждого этапа составьте свою таблицу. Равенства времени определяются отдельно для каждого этапа.

\subsection*{Пример: автомобиль и мотоциклист}

Расстояние между $A$ и $B$ составляет $80$ км. Автомобиль выехал из $A$ первым, а через $20$ минут из $A$ вслед за ним выехал мотоциклист со скоростью $90$ км/ч. Мотоциклист догнал автомобиль в точке $C$, повернул обратно и проехал половину расстояния от $A$ до $C$ за то время, пока автомобиль доехал от $C$ до $B$.

Обозначим $AC=x$ км, а скорость автомобиля через $y$ км/ч. Тогда $CB=80-x$, где $0<x<80$ и $0<y<90$.

\textbf{Этап I: от старта до встречи в $C$.} У обоих путь $x$, но автомобиль стартовал на $20$ минут раньше.
\begin{center}
\renewcommand{\arraystretch}{1.32}
\begin{tabularx}{\linewidth}{@{}lXXX@{}}
\toprule
Участник & $v$, км/ч & $t$, ч & $S$, км\\
\midrule
Автомобиль & $y$ & $x/y$ & $x$\\
Мотоциклист & $90$ & $x/90$ & $x$\\
\bottomrule
\end{tabularx}
\end{center}
Отсюда $\displaystyle\frac{x}{y}-\frac{x}{90}=\frac{20}{60}=\frac13$.

\textbf{Этап II: после встречи в $C$.} Автомобиль проезжает $80-x$ км, мотоциклист возвращается на $x/2$ км. Теперь они начинают этап \emph{одновременно}, поэтому времена равны.
\begin{center}
\renewcommand{\arraystretch}{1.32}
\begin{tabularx}{\linewidth}{@{}lXXX@{}}
\toprule
Участник & $v$, км/ч & $t$, ч & $S$, км\\
\midrule
Автомобиль & $y$ & $(80-x)/y$ & $80-x$\\
Мотоциклист & $90$ & $x/180$ & $x/2$\\
\bottomrule
\end{tabularx}
\end{center}
Отсюда $\displaystyle\frac{80-x}{y}=\frac{x}{180}$.

Вместе получаем систему:
\[
\boxed{\begin{cases}
\dfrac{x}{y}-\dfrac{x}{90}=\dfrac13,\\[5pt]
\dfrac{80-x}{y}=\dfrac{x}{180}.
\end{cases}}
\]
\key{Для данного этапа подготовки} важно безошибочно составить обе таблицы и систему. Решение таких систем будет отработано отдельно.

\clearpage
\section*{Алгоритм: как составить уравнение по задаче на движение}
\vspace{6mm}
\begin{center}
\begin{tikzpicture}[
  >={Latex[length=2.4mm]},
  start/.style={ellipse,draw=darkaccent,thick,minimum width=4cm,minimum height=1cm,align=center,fill=white},
  act/.style={rounded corners=3pt,draw=darkaccent,thick,text width=10.8cm,minimum height=1.15cm,align=center,inner sep=6pt,fill=white},
  dec/.style={diamond,aspect=2.75,draw=accent,thick,text width=4.1cm,align=center,inner sep=2pt,fill=white},
  choice/.style={rounded corners=3pt,draw=accent,thick,text width=4.0cm,minimum height=1.3cm,align=center,inner sep=5pt,fill=white},
  arrow/.style={->,thick,draw=darkaccent}
]
\node[start] (start) at (0,0) {Начало};
\node[act] (stage) at (0,-1.75) {Выделите участников и этапы: до встречи, после встречи, до финиша};
\node[act] (scheme) at (0,-3.6) {Нарисуйте схему: направления, расстояния, задержки и точки встречи};
\node[act] (table) at (0,-5.45) {Заполните таблицу $v$--$t$--$S$ для каждого этапа, используя $S=vt$};
\node[dec] (decision) at (0,-8.1) {Что сравниваем?};
\node[choice] (path) at (-3.5,-11.1) {Пути:\ $S_{\text{быстр}}-S_{\text{медл}}=d$};
\node[choice] (time) at (3.5,-11.1) {Времена:\ $t_{\text{раньше}}-t_{\text{позже}}=\Delta t$};
\node[act] (equation) at (0,-13.1) {Составьте уравнение или систему. Переведите минуты в часы и метры в километры};
\node[start] (end) at (0,-15.2) {Проверьте смысл и ограничения};
\draw[arrow] (start)--(stage);
\draw[arrow] (stage)--(scheme);
\draw[arrow] (scheme)--(table);
\draw[arrow] (table)--(decision);
\draw[arrow] (decision.south west) -| (path.north);
\draw[arrow] (decision.south east) -| (time.north);
\draw[arrow] (path.south) -- ++(0,-.48) -| ([xshift=-2cm]equation.north);
\draw[arrow] (time.south) -- ++(0,-.48) -| ([xshift=2cm]equation.north);
\draw[arrow] (equation)--(end);
\end{tikzpicture}
\end{center}
\clearpage

\section*{7. Ошибки, которые нужно предупреждать}
\begin{enumerate}
\item \textbf{Смешение единиц.} Если $v$ выражена в км/ч, то $300$ м записывают как $0{,}3$ км, а $20$ минут --- как $\frac13$ часа.
\item \textbf{Неправильная разность.} Если один ехал дольше, вычитайте из \emph{большего} времени \emph{меньшее}, чтобы получить положительную задержку.
\item \textbf{Путаница с направлением.} После разворота пройденный путь не равен расстоянию от старта до точки встречи. Путь быстрого участника в примере с опушкой равен $6+x$.
\item \textbf{Лишняя переменная пугает.} В задаче на догонку скорость $u$ может сократиться. Не нужно заранее вводить ещё одно уравнение, если искомое уже определяется.
\item \textbf{Одно уравнение на два разных этапа.} Если направление или условия изменились, заполняйте вторую таблицу и составляйте отдельное соотношение.
\item \textbf{Деление на нуль.} В дробях $S/v$ скорость в знаменателе должна быть положительна. Например, при скоростях $x$ и $x-12$ допустимо только $x>12$.
\end{enumerate}

\section*{Итоговая памятка}
\begin{center}
\renewcommand{\arraystretch}{1.34}
\begin{tabularx}{\linewidth}{@{}p{.40\linewidth}X@{}}
\toprule
Сигнал из текста & Какое равенство искать\\
\midrule
«Догнал», «на столько-то впереди» & Разность пройденных путей равна исходному отрыву.\\
«Между ними стало $d$» & Разность путей равна $d$.\\
«Выехал позже, прибыл вместе» & Разность полных времён движения равна задержке.\\
«Дошёл до точки и вернулся» & Учитывайте оба участка пути и сравнивайте полные времена.\\
«После встречи...» & Начните второй этап: новая таблица и новое равенство времён или путей.\\
\bottomrule
\end{tabularx}
\end{center}

\clearpage
\section*{Тренировочный блок --- путь А}
\smallnote{10 задач. Выполняйте задания последовательно: от прямой догонки к движению с задержкой старта и двумя этапами. В задачах 9--10 требуется составить таблицы и систему; решать систему пока не нужно.}
\vspace{3mm}

\subsection*{Базовый уровень}
\task{1}{Догонка с первоначальным отрывом}
По одной прямой дороге в одном направлении одновременно начали движение легковой автомобиль и грузовик. Легковой автомобиль выехал из пункта $A$, а грузовик --- из расположенного на $96$ км впереди пункта $B$. Скорость автомобиля составляет $64$ км/ч, скорость грузовика --- $48$ км/ч. Оба продолжают движение без остановок. Через сколько часов после начала движения легковой автомобиль догонит грузовик? Нарисуйте схему и составьте таблицу движения.

\task{2}{Пешеходы из одной точки}
Два туриста одновременно вышли из одного парка и пошли по одной аллее в одном направлении. Первый двигался с постоянной скоростью $6$ км/ч, второй --- со скоростью $4$ км/ч. Через некоторое время расстояние между ними стало равно $500$ м. Через сколько минут это произошло? Запишите разность пройденных путей, предварительно согласовав единицы измерения.

\task{3}{Расстояние между велосипедистами}
Два велосипедиста одновременно выехали от одной станции по одной дороге в одном направлении. Скорость первого составляет $18$ км/ч, скорость второго --- $16$ км/ч. Каждый едет равномерно и без остановок. Какое расстояние в километрах окажется между велосипедистами через $45$ минут после старта? Обоснуйте ответ через разность пройденных расстояний, а не через отдельное сложение путей.

\task{4}{Неизвестная скорость грузовика}
По прямой дороге в сторону города $C$ одновременно начали движение легковой автомобиль и грузовик. Легковой автомобиль выехал из города $A$, а грузовик --- из пункта $B$, расположенного на $150$ км дальше по той же дороге. Скорость легкового автомобиля равна $95$ км/ч, а скорость грузовика неизвестна. Известно, что автомобиль догнал грузовик ровно через $5$ часов после выезда. Найдите скорость грузовика и проверьте, что она меньше скорости автомобиля.

\Needspace{15\baselineskip}\subsection*{Средний уровень}
\task{5}{Поздний старт легкового автомобиля}
Из посёлка по прямому шоссе выехал грузовик со скоростью $50$ км/ч. Ровно через час из того же посёлка и в том же направлении выехал легковой автомобиль со скоростью $75$ км/ч. Оба транспортных средства ехали без остановок и не меняли скорость. Через сколько часов \emph{после своего выезда} легковой автомобиль догонит грузовик? Сколько часов к этому моменту будет находиться в пути грузовик? Покажите, почему их времена движения до догонки не равны.

\task{6}{Пешеход дошёл до разворота}
Два туриста одновременно вышли из одного места и направились к смотровой площадке, расположенной в $4$ км от старта. Первый турист идёт со скоростью $3$ км/ч, второй --- со скоростью $5$ км/ч. Быстрый турист дошёл до площадки и сразу повернул обратно, продолжая идти с той же скоростью. На обратном пути он встретил первого туриста. На каком расстоянии от смотровой площадки произошла встреча? Используйте переменную для расстояния от площадки и сравните полные времена обоих туристов.

\task{7}{Одинаковое расстояние, разные моменты старта}
Грузовик и автобус проезжают от станции $A$ до станции $B$ ровно по $60$ км каждый. Сначала выехал грузовик, а через $15$ минут вслед за ним отправился автобус. Автобус двигался на $20$ км/ч быстрее грузовика; оба прибыли на станцию $B$ одновременно. Найдите скорости грузовика и автобуса. Начните с обозначения скорости грузовика через $x$ км/ч, запишите времена движения дробями и укажите допустимые значения $x$.

\clearpage
\Needspace{15\baselineskip}\subsection*{Повышенный уровень}
\task{8}{Поезда с задержкой отправления}
С одной станции по одному пути в одном направлении отправился грузовой поезд со скоростью $54$ км/ч. Через $40$ минут после него по тому же пути отправился скорый поезд, скорость которого составляет $90$ км/ч. Считайте, что оба поезда идут равномерно, без остановок и без изменения скорости. На каком расстоянии от станции скорый поезд догонит грузовой? При составлении уравнения различайте время движения грузового и скорого поездов.

\task{9}{Догонка и разворот мотоциклиста: две таблицы}
Расстояние между городами $A$ и $B$ по прямой дороге равно $120$ км. Из города $A$ в город $B$ выехал автомобиль, а через $15$ минут вслед за ним из того же города отправился мотоциклист со скоростью $80$ км/ч. Мотоциклист догнал автомобиль в промежуточной точке $C$, сразу повернул обратно и за время, пока автомобиль доехал от $C$ до $B$, проехал в обратном направлении треть расстояния от $A$ до $C$. Обозначьте расстояние $AC$ через $x$ км, а скорость автомобиля --- через $y$ км/ч. Изобразите движение на схеме, составьте \emph{две таблицы} (до встречи и после неё) и получите систему двух уравнений. \textbf{Решать систему не требуется.}

\task{10}{Догонка и общий финиш: два этапа}
Из пункта $A$ по прямой дороге в пункт $B$, расстояние между которыми равно $270$ км, выехал грузовик. Через полтора часа вслед за ним из пункта $A$ выехал легковой автомобиль, который ехал быстрее грузовика. Спустя ровно $3$ часа собственного движения автомобиль догнал грузовик, после чего оба продолжили путь до $B$ с прежними скоростями. Легковой автомобиль прибыл в $B$ на $45$ минут раньше грузовика. Обозначьте скорость грузовика через $v$ км/ч, а скорость автомобиля --- через $w$ км/ч. Составьте отдельные таблицы для движения \emph{до догонки} и \emph{после догонки} и запишите два уравнения, связывающие $v$ и $w$. \textbf{Решать систему не требуется.}

\clearpage
\section*{Ответы и ориентиры для самопроверки}
\smallnote{В заданиях 9--10 проверяется правильность математической модели, а не решение получившейся системы.}
\vspace{4mm}
\renewcommand{\arraystretch}{1.5}
\begin{longtable}{@{}p{.075\linewidth}p{.855\linewidth}@{}}
\toprule
№ & Ответ или проверочное уравнение\\
\midrule
\endhead
1 & $t=\dfrac{96}{64-48}=\boxed{6\ \text{ч}}$.\\
2 & $500\ \text{м}=0{,}5\ \text{км}$; $\ (6-4)t=0{,}5$, поэтому $\boxed{15\ \text{мин}}$.\\
3 & $t=\dfrac{45}{60}=\dfrac34\ \text{ч}$; $\ (18-16)t=\boxed{1{,}5\ \text{км}}$.\\
4 & $(95-v)\cdot5=150$; $\boxed{v=65\ \text{км/ч}}$.\\
5 & $75t=50(t+1)$; автомобиль ехал $\boxed{2\ \text{ч}}$, грузовик --- $\boxed{3\ \text{ч}}$.\\
6 & $\dfrac{4-x}{3}=\dfrac{4+x}{5}$; $\boxed{x=1\ \text{км}}$ от площадки.\\
7 & $\dfrac{60}{x}-\dfrac{60}{x+20}=\dfrac14$, $x>0$; грузовик --- $\boxed{60\ \text{км/ч}}$, автобус --- $\boxed{80\ \text{км/ч}}$.\\
8 & $54\left(t+\dfrac23\right)=90t$; скорый поезд ехал $1$ час, расстояние --- $\boxed{90\ \text{км}}$.\\
9 & При $0<x<120$, $0<y<80$:
$\displaystyle\begin{cases}
\frac{x}{y}-\frac{x}{80}=\frac14,\\[4pt]
\frac{120-x}{y}=\frac{x}{240}.
\end{cases}$
Второе уравнение относится \emph{только к движению после встречи}.\\
10 & До догонки грузовик ехал $4{,}5$ ч, автомобиль --- $3$ ч. После догонки остаётся $270-4{,}5v=270-3w$ км:
$\displaystyle\begin{cases}
4{,}5v=3w,\\[4pt]
\frac{270-4{,}5v}{v}-\frac{270-3w}{w}=\frac34.
\end{cases}$
Ограничения: $0<v<w$, $4{,}5v<270$.\\
\bottomrule
\end{longtable}

\vfill
\begin{center}
\small\color{muted} Лёвин Артём Александрович --- эксклюзивно для Ярослава Гаврилова
\end{center}
\end{document}
